\chapter{Aufnahme eines Computertomogramms} \label{sec:ct_chapter}

\section{Funktionsprinzip der Computertomographie}

Es soll in diesem Versuchsteil das Prinzip der Computertomographie (CT), eine wichtige medizinische Diagnosetechnik \cite{medizinphysik}, demonstriert werden \cite[2]{529va}.
\todo{prinzip kurz erklären, was ist eine projektion und was eine rekonstruktion, was heißen die pixelstufen dann am ende?}

\section{Aufbau \& Durchführung}

Es wurde zunächst die Wolfram (W)-Röhre in das Röntgengerät eingebaut. Diese Röhre wurde anstatt der Mo-Röhre verwendet, da in der Gebrauchsanweisung des CT-Apparats beide Röhren gelistet sind, wobei die W-Röhre aber empfohlen wird, da sie eine höhere Intensität des Bilds auf dem Floreszenzschirms bei gleichen Betriebsparametern liefert \cite[2]{LDctmodul}. Es war zu erwarten, dass das Bild so besser erkennbar wird. Es wurde alles bis auf das Goniometer aus dem Experiementierraum entfernt und das Goniometer ganz nach rechts geschoben, damit das scanbare Volumen möglichst groß wird \cite[3]{LDctmodul}. Als zu untersuchendes Objekt wurde jetzt auf einem Targethalter ein von Styropor umgebener getrockneter kleiner Frosch in das Röntgengerät mithilfe der Halterung am Goniometer eingebaut.

Als nächster Schritt wurde der Aufbau korrekt justiert. Hierzu wurde die Justageanleitung der Gebrauchsanweisung befolgt \cite{LDctmodul}.
\todo{justage}

\section{Durchführung}

\section{Auswertung}

\jafpp{frosch_ct/frog_greyscale}{}{frosch_ct/frog_greyscale_cropped}{}
\jafpp{frosch_ct/frog_sideways}{}{frosch_ct/frog_sideways_cropped}{}
\jafpp{frosch_ct/frog_hohe_transparenz}{}{frosch_ct/frog_niedr_transparenz}{}
\jafpp{frosch_ct/frog_upsidedown}{}{frosch_ct/frog_querdurch}{}
\jafps{frosch_ct/frosch}{}{0.45}
